\section{Related Work}
\label{sec:related}

\textbf{Device fingerprinting.} Devices can be identified from what they emit, even without any
access to the payload. Kohno et al.\ fingerprint a remote host from the skew of its clock as seen
in TCP timestamps~\cite{kohnoRemotePhysicalDevice2005}, and GTID fingerprints a device and its
type from the timing signature of its packets~\cite{radhakrishnanGTIDTechniquePhysical2015}.
These works define the threat we address; our objective, on the other hand, is to remove the timing
features they rely on from the master-facing observation. Their target is the device type, whereas
the classifier we evaluate separates transaction classes on one outstation, so we suppress a feature
their cross-layer response-time method uses without demonstrating that two devices become
indistinguishable.

\textbf{ICS and power-system fingerprinting.} Formby et al.\ brought fingerprinting to control
systems with the cross-layer response time and the physical operation time, measured at the
device, of DNP3 outstations~\cite{formbyWhosControlYour2016}, and Gu et al.\ added device physics
as a feature~\cite{guFingerprintingCyberPhysicalSystem2018}. Jeon et al.\ fingerprint SCADA
devices passively without deep packet inspection~\cite{jeonPassiveFingerprintingSCADA2016}. The
reconnaissance these techniques serve enables false-data injection and process-control
attacks~\cite{liuFalseDataInjection2009, cardenasAttacksProcessControl2011}. Defenses in this
setting have so far worked at the content and connectivity layer: DefRec virtualizes physical
functions to present deceptive content~\cite{linDefRecEstablishingPhysical2020}, RAINCOAT
randomizes data acquisition and injects decoy
measurements~\cite{linRAINCOATRandomizationNetwork2019}, and a companion testbed evaluates such
anti-reconnaissance approaches on grid infrastructure~\cite{linCyberPhysicalTestbedCase2020}.
Compared to these, our framework works below the semantics they reshape and complements them: it
changes when the wire carries the outstation's answer, not what the answer says.

\textbf{General traffic obfuscation.} Traffic morphing transforms one class's packet-size
distribution into another's~\cite{wrightTrafficMorphingEfficient2009}, while Dyer et al.\ show
that per-packet padding and fragmentation still leave exploitable
features~\cite{dyerPeekaBooStillSee2012}. Tamaraw fixes packet size and rate at a bandwidth
cost~\cite{caiSystematicApproachDeveloping2014}, and WTF-PAD pads inter-packet gaps
adaptively~\cite{juarezEfficientWebsiteFingerprinting2016}. However, deep fingerprinting attacks
defeat several of these~\cite{sirinamDeepFingerprintingUndermining2018}, and their accuracy at
Internet scale is contested~\cite{panchenkoWebsiteFingerprintingInternet2016}. The same
size-and-timing channel identifies smart-home devices under encryption, where shaping is the
countermeasure~\cite{apthorpeKeepingSmartHome2019}, and Pacer and NetShaper shape a tenant's
traffic to a schedule that bounds side-channel leakage in the
cloud~\cite{mehtaPacerComprehensiveNetwork2022, sabziNetShaperDifferentiallyPrivate2024}. These
defenses assume an encrypted payload and a cooperating end host or hypervisor that can add cover
traffic. Our framework, in contrast, targets a plaintext protocol whose feature is one interval
inside a request-response exchange, adds no byte, and runs where the outstation cannot be
changed.

\textbf{Programmable data-plane traffic transformation.} P4 defines the programmable parser and
match-action model we use~\cite{bosshartP4ProgrammingProtocolIndependent2014}. PIFO defines
programmable scheduling at line rate~\cite{sivaramanProgrammablePacketScheduling2016}, and SP-PIFO
approximates it with strict-priority queues on today's
switches~\cite{alcozSPPIFOApproximatingPushIn2020}, which is the primitive our reservoirs use.
Ditto pads and injects chaff to a fixed pattern at line rate~\cite{meierDittoWANTraffic2022},
Minos morphs and schedules encrypted traffic on a programmable switch with low
overhead~\cite{wangMinosLightweightDynamic2025}, and Securitas mixes fragmentation and insertion
under a learned policy across several data planes~\cite{xieSecuritasDefendingTraffic2026}.
NetHide obfuscates topology in the data plane against link-flooding
reconnaissance~\cite{meierNetHideSecurePractical2018}. Closest to our setting, Hu et al.\ use
programmable switches to enhance industrial-protocol
security~\cite{huIndustrialNetworkProtocol2023}. Our framework differs from these in the
observable it addresses, i.e., the master-visible timing of a DNP3 transaction, and in holding
packets rather than adding or reshaping them.

\textbf{DNP3 timing and security.} East et al.\ catalogue attacks on DNP3 and the absence of
encryption that makes its semantics visible~\cite{eastTaxonomyAttacksDNP32009}, and
specification-based and state-based intrusion detection has been adapted to
DNP3~\cite{linAdaptingBroSCADA2013, fovinoModbusDNP3StateBased2010}. These do not change the
timing an observer sees. On the control path we report a different quantity from theirs: the master-visible OPERATE
response-to-acknowledgment interval, which is a protocol-response timing feature. Formby et al.\
estimate physical operation time from a sequence-of-events-recorder timestamp carried in the
outstation's application payload, and report that the alternative estimate from unsolicited-response
arrival times produced no usable results~\cite{formbyWhosControlYour2016}. A mechanism that moves
packets in time and changes no byte cannot alter such a payload timestamp, and the traffic we
evaluate carries neither unsolicited responses nor any time-bearing object. We therefore make no
claim about physical operation time, and we do not present our control-path result as a correlate of
theirs.
